% !TeX root = ../papers/nc-p3.tex
\section{十五次元成分の全点列に対する行列恒等式と分岐周期の上界}
\label{sec:pi-periodic-point-paths}
$B$ は標準四頂点で折り畳んだ代数とする。点列の長さは文字の個数で数え、
周期は頂点ラベルも含めて最小に取る。$s_m$ は $m$ 変数の積の交代和である。
\begin{theorem}\label{thm:pi-periodic-point-paths}
十五次元成分の閉包の任意の標準AS正則窓について次が成立する。
\begin{enumerate}
\item 周期点列の最小周期は $4,8,12,16$ のいずれかである。
順または逆の分岐を一度でも訪れる周期列の最小周期は $4,8,12$ に限る。
従って曲面射影を持たない分岐窓では、定理\ref{thm:periodic-branch-obstruction}の
分岐を通る周期列の長さの上界を40から12にできる。
\item 左右点ファイバーなら、グラフの次元を問わず、全閉点が周期点である。
四移動 $\sigma$ の各閉点での周期は $1,2,3,4$ のいずれかであり、
$X_{\mathrm{red}}$ 上では $\sigma^{12}=1$ である。
\item 二つの周期列を、途中で $\ell\ge2$ 個の文字を通って切り替えられるとする。
すなわち最初の列のある状態からその後続文字と異なる文字で出て、
$\ell$ 番目では第二列の到着状態への周期からの入射文字と異なる文字で入る。
全隣接対が点対応に属することを要求する。それぞれの最小周期を $p,q$ とすると
\[
                 p+q+\min(\ell-1,4)\le16 .
\]
とくに順序を除く周期の組は $(4,4)$ または $(4,8)$ に限る。
\end{enumerate}
\end{theorem}
\begin{proof}[全特殊点に残る恒等式]
系\ref{cor:quartic-division-corners}により成分の一般折り畳み代数は次数16の中心単純環に入り、
各隅は次数4の可除代数に入る。分裂体でのAmitsur--Levitzki恒等式から
$B$ は $s_{32}$、各隅は $s_8$ を満たす。標準正則な特殊点には、この成分の一般点から
離散付値環を通じて至ることができる。一般・特殊の全経路次元が標準値なので
各次数の経路格子は自由である。したがって任意の有限個の経路の線型結合に恒等式を代入し、
分数体上での零を格子上の零として還元すれば、両恒等式が特殊点でも成立する。
特殊点の隅が整域であるという条件は使わない。

周期 $p$ の点列から、一周の各状態に一元を置き、文字をその点で評価して次の状態に移る
有限次元表現を作る。その像は $M_p$ 全体である。実際、同じ頂点で始まる異なる状態の
一周の評価テンソルは射影的に異なる（等しければより短い周期になる）。
各相手の一周を消して指定状態の一周で非零となる線型結合を選び、
相手の個数だけ一周の区間を繰り返して掛けると、指定状態だけに作用する非零対角単位が得られる。
隣状態どうしは非零文字で結べるので全行列単位が得られる。
$m$ 個の直交冪等元 $e_1,\ldots,e_m$ と各 $e_jAe_{j+1}$ の元で
全積が非零なものがあれば、$s_{2m-2}$ は消えない。
実際 $e_1,a_1,e_2,a_2,\ldots,e_{m-1},a_{m-1}$ を代入するとこの順の積だけが残る。
従って $p\le16$ である。

周期列の順分岐では、元の文字と比例しない別の後続文字を使い、
その一歩だけを新しい終端の一次元状態に着ける。
追加した長さ二の経路は一つ前の文字とこの後続文字を通るので、これは依然 $B$ の表現である。
行ベクトルの右作用で終端は頂点単純、商は前の周期単純である。この拡大は分裂しない。
実際、周期状態の持上げを終端の元の定数倍だけ変えたとき、
追加の係数を消すにはそれが周期の同じ位置の文字の倍数であることが必要になる。
像の二つの対角表現は非同型単純でそれぞれ全行列環なので、半単純商は $M_p\oplus k$ である。
非分裂性から非対角根基は非零で、単純な矩形行列両側加群なのでその全体となる。
従って像は二ブロック上三角環であり、上の行列単位列を $p+1$ 個まで取れる。
$p+1\le16$ と四の倍数性から $p\le12$ を得る。逆側は反対環で同じである。
\end{proof}
\begin{proof}[決定的な点列の非周期性の排除]
グラフの四移動で非周期な閉点があれば、その軌道の相異なる点は永久に一致しない。
連続する64文字を八文字ずつ八区間に切る。各区間の相対位置を一つ固定すれば、
そこで八区間に現れる点は同じ文字空間内で相異なる。
目標の一区間を残すため、七つの他区間それぞれに異なる相対位置を割り当て、
そこで他区間の点を消し目標の点で非零な文字を取る。余った一位置も目標で非零に取れば、
この長さ8の経路は目標だけで非零である。八区間全てを順に目標にして規格化し、
長さ8の八経路を得る。いずれも同じ四頂点の隅に属する。これらの $s_8$ を
元の64文字の点列で評価すると、八区間の評価行列は単位行列なので行列式は1である。
隅の恒等式に反する。よって周期的で、第一の上界から四移動十二回が全閉点を固定する。
被約有限型スキーム上の射は閉点上で決まる。
\end{proof}
\begin{proof}[二周期を結ぶ場合]
二つの周期の状態を別々に取り、最初の状態 $A_s$ から最後の状態 $B_0$ への
$\ell$ 本の矢と途中の $\ell-1$ 個の一次元状態を追加する。
二次関係は各二歩の係数が点対応で消えるため全て成立する。
各頂点から出る一文字ずつを全ての評価で非零に選んで和を $x$ とする。
各周期の各矢の代表倍率を選び、一方と他方で $x$ の一周の積の倍率を独立に調節できる。
$p,q,4$ の公倍数 $N\ge\ell$ を取り、$x^N$ の両周期上の倍率を相異なる非零値にする。
中間ブロックは冪零である。$x^N$ の三つの一般固有空間への射影 $E_A,E_B,E_T$
はその多項式で表され、元の四頂点冪等元と可換である。
各周期への圧縮環は、射影像の次元と周期商上の行列単位から $M_p,M_q$ 全体である。

最初の追加文字で非零、同じ状態での周期文字で零となる元を第一の矢に、
最後の追加文字で非零、$B_0$ へ入る周期文字で零となる元を最後の矢に使う。
中間の各矢では所定の点で非零な元を使う。最初の $k=\min(\ell-1,4)$ 個の途中状態の直後に、
その頂点冪等元と $E_T$ の積を挿入する。この $k$ 個は位置を法4で異にするから互いに直交する。
この積は $E_A$ と $E_B$ の間で依然非零である。実際 $A_s E_A$ は $A_s$ と
中間・終周期の補正だけからなる。第一矢で始周期に残る項は消え、中間の第一状態への係数は非零となる。
$E_T$ は中間ブロックで恒等であり、中間状態からの補正は終周期にだけ出る。
従って各段で所定の途中状態への係数を保ち、他の項はそれより先の状態か終周期にある。
最後の矢では終周期からの $B_0$ への係数は零なので、追加経路の非零係数が打ち消されない。
終周期では $E_B$ は恒等である。

両端を各圧縮環の行列単位で挟み、非零な一成分を選ぶ。
その前後にそれぞれの全対角単位を通る行列単位列をつなげば、
$p+q+k$ 個の相異なる直交冪等元を通る非零積になる。
17個以上ならその先頭17個で $s_{32}$ が非零となって矛盾する。以上で上界を得る。
\end{proof}

\begin{corollary}\label{cor:periodic-branch-central-basepoints}
任意の標準正則窓で、周期列の右点加群上に正次数の斉次中心元が非零に作用するとする。
その次数は列の最小周期の倍数であり、この列上には順分岐も逆分岐もない。
特に四次中心の一つでも非零に作用する周期列は長さ4で分岐がない。
分岐を訪れる全周期列は全正次数中心元の共通零点に限られる。
\end{corollary}
\begin{proof}
折り畳みの各頂点に成分を持つ中心元 $z$ の次数を $d$ とし、点加群の各次数での作用係数を $c_j$ とする。
矢との可換式は $c_j p_{j+d}=c_{j+1}p_j$ である。どの $p$ も非零なので、
一つの $c_j$ が非零なら全て非零で、列が $d$ 移動で戻る。右点加群は $k[z]$ 上で
最初の $d$ 個の同次元を基底とする自由加群である。従って右 $B$ 加群として有限表示である。
実際この基底各元と有限個の代数生成元の積を基底の $k[z]$ 線型結合で記すと、
これら有限個の関係で任意の語を簡約でき、簡約後の独立性は $k[z]$ 上の基底性から従う。
しかし順中心を繰り返せば定理\ref{thm:periodic-branch-obstruction}の第一Tor計算は
無限個の次数で新しい関係を与える。この計算は曲面の有無を問わない。
逆順で同じ周期の左点加群を取ると、$z$ の経路評価も同じ（開始位置を取り直すだけ）なので
左側の同じ議論で逆分岐も排除できる。
\end{proof}
